\chapter{Formulação QUBO para Escalonamento}
\label{chap:formulacao-qubo-escalonamento}

\section{CPU Scheduling em Sistemas Operativos}
\label{sec:cpu-scheduling}

Um escalonador de \acs{CPU} determina que processo ou \textit{thread} executa em cada núcleo do processador
em cada instante. Escalonadores como o \ac{CFS} do Linux consideram
simultaneamente latência, \textit{throughput}, prioridade, afinidade de núcleo e localização de cache,
tomando decisões em microssegundos com base em heurísticas clássicas~\cite{LinuxCFS}.

No entanto, a decisão de atribuição de processos a núcleos é fundamentalmente um problema de partição combinatorial. Dado um conjunto de $N$ processos com cargas $w_i$ e $K$ núcleos, o número de atribuições possíveis cresce como $K^N$, tornando a procura exaustiva para instâncias reais~\cite{Kochenberger2014QUBO,Glover2019QUBO}.

Este projeto adota uma formulação simplificada deste problema: cada \textit{workload entity} possui
uma utilização de \acs{CPU} conhecida, e o objetivo é distribuí-las pelos núcleos disponíveis minimizando o
desequilíbrio de carga. Esta abstração isola a componente combinatorial do escalonamento sem replicar a complexidade de um escalonador de kernel completo.

\section{Formulação QUBO}
\label{sec:formulacao-qubo}

Um problema \acs{QUBO}~\cite{Kochenberger2014QUBO,Glover2019QUBO} consiste em minimizar
uma expressão quadrática sobre variáveis binárias $x_i \in \{0, 1\}$, representada de forma compacta como:

\begin{equation}
    \min_{\mathbf{x} \in \{0,1\}^n} \mathbf{x}^\top Q\, \mathbf{x} = \sum_{i} \sum_{j} Q_{ij} \cdot x_i \cdot x_j
\end{equation}
\listeq{Forma geral de um problema QUBO}

onde $Q \in \mathbb{R}^{n \times n}$ é uma matriz cujos coeficientes codificam o problema a resolver.
Expandindo o somatório, surgem dois tipos de termos:

\begin{itemize}
    \item \textbf{Termos diagonais} ($i = j$): como $x_i^2 = x_i$ para variáveis binárias,
    $Q_{ii} \cdot x_i \cdot x_i = Q_{ii} \cdot x_i$. Estes termos associam um custo individual
    a cada variável --- um coeficiente $Q_{ii} < 0$ incentiva o otimizador quântico a escolher $x_i = 1$; um coeficiente
    $Q_{ii} > 0$ incentiva $x_i = 0$.

    \item \textbf{Termos fora da diagonal} ($i \neq j$): $Q_{ij} \cdot x_i \cdot x_j$ contribui
    para a energia apenas quando $x_i = x_j = 1$. Um coeficiente positivo penaliza a ativação
    simultânea de $x_i$ e $x_j$; um coeficiente negativo recompensa-a.
\end{itemize}

A energia de uma solução $\mathbf{x}$ é o valor $\mathbf{x}^\top Q\, \mathbf{x}$.
A solução ótima $\mathbf{x}^*$ é a atribuição binária que minimiza esta energia no espaço de
$2^n$ configurações possíveis.

\section{Construção da Matriz Q}
\label{sec:construcao-matriz-q}

Dado um conjunto de $N$ processos com cargas $w_i$ e $K$ núcleos, define-se uma
variável binária $x_{ij} \in \{0,1\}$ para cada par processo-núcleo. Esta variável
toma o valor $1$ quando o processo $i$ é atribuído ao núcleo $j$, e o valor $0$
caso contrário.
Na implementação, a variável $x_{ij}$ é armazenada no índice linear $iK+j$ da matriz $Q$.

O objetivo original do problema é equilibrar a carga entre núcleos. Para isso,
mede-se o desequilíbrio através do desvio quadrático total das cargas relativamente
à carga média:

\begin{equation}
    \min \sum_{j=0}^{K-1}(L_j - L_{avg})^2.
\end{equation}
\listeq{Desvio quadrático total das cargas}

onde
\begin{equation}
    L_{avg} = \frac{W_{total}}{K}.
\end{equation}
\listeq{Carga média alvo}

Para obter a formulação \acs{QUBO}, este termo é expandido e combinado com penalidades que impõem
atribuições válidas. O resultado é codificado na matriz $Q$, de modo que o \textit{solver}
minimize a energia
\begin{equation}
    E(\mathbf{x}) = \mathbf{x}^{\top}Q\mathbf{x}.
\end{equation}
\listeq{Energia QUBO de uma solução}

Assim, minimizar a energia \acs{QUBO} corresponde, por construção, a procurar uma atribuição
simultaneamente equilibrada e válida.

Para cada núcleo $j$, a carga total é dada por:
\begin{equation}
    L_j = \phi_j + \sum_{i=0}^{N-1} w_i x_{ij},
\end{equation}
\listeq{Carga total no núcleo \(j\)}
onde $\phi_j$ representa carga já comprometida nesse núcleo, relevante para a secção da
decomposição iterativa. Na primeira iteração, ou num \acs{QUBO} sem carga fixa, tem-se
$\phi_j = 0$.

Substituindo $L_j$ no termo de um núcleo e usando $x_{ij}^2=x_{ij}$, obtém-se:
\begin{equation}
\begin{aligned}
(L_j - L_{avg})^2
&= \left(\phi_j + \sum_{i=0}^{N-1}w_ix_{ij} - L_{avg}\right)^2 \\
&= \underbrace{\sum_{i=0}^{N-1} d_{ij}x_{ij}}_{\text{termos diagonais}}
 + 2\underbrace{\sum_{0 \leq i < k}^{N-1} w_iw_kx_{ij}x_{kj}}_{\text{termos de co-localização}}
 + \underbrace{(\phi_j - L_{avg})^2}_{\text{constante}},
\end{aligned}
\end{equation}
\listeq{Expansão do termo de balanceamento}
onde
\begin{equation}
    d_{ij}=w_i^2 + 2w_i\phi_j - 2L_{avg}w_i.
\end{equation}
\listeq{Coeficiente diagonal do balanceamento}

Estes coeficientes são então
adicionados à matriz $Q$. Os termos diagonais estão associados
à ativação individual de uma variável $x_{ij}$; os termos de co-localização aparecem
quando dois processos distintos são atribuídos ao mesmo núcleo.

Além do objetivo de balanceamento, é necessário garantir que cada processo é atribuído a
exatamente um núcleo. Para isso adiciona-se a penalidade one-hot:
\begin{equation}
    P \sum_{i=0}^{N-1}\left(\sum_{j=0}^{K-1}x_{ij} - 1\right)^2.
\end{equation}
\listeq{Penalidade one-hot}
Esta penalidade segue a prática comum de converter restrições de igualdade em termos
quadráticos penalizados em \acs{QUBO}~\cite{Glover2019QUBO}. Neste caso, contribui com $-P$
na diagonal de cada variável e com $+P$ entre variáveis do mesmo processo em núcleos diferentes.

Combinando o objetivo de balanceamento e a penalidade one-hot, cada entrada da matriz
é determinada pelo tipo de relação entre as duas variáveis consideradas:
\begin{equation}
Q[x_{ij},x_{kl}] =
\begin{cases}
w_i^2 + 2w_i\phi_j - 2L_{avg}w_i - P, & \text{se } i=k \text{ e } j=l,\\
w_iw_k, & \text{se } i \neq k \text{ e } j=l,\\
P, & \text{se } i=k \text{ e } j \neq l,\\
0, & \text{se } i \neq k \text{ e } j \neq l.
\end{cases}
\end{equation}
\listeq{Coeficientes da matriz \(Q\)}
A primeira linha corresponde à contribuição individual de cada variável. A segunda
linha penaliza processos distintos que partilham o mesmo núcleo, enquanto a terceira
corresponde à penalidade one-hot entre núcleos diferentes do mesmo processo. A última
linha indica que variáveis de processos e núcleos diferentes não interagem diretamente.

\section{Exemplo Ilustrativo}
\label{sec:exemplo-qubo}

Considere-se $N=2$ processos e $K=2$ núcleos, com cargas $w_0=0.3$ e $w_1=0.5$,
sem carga fixa ($\phi_j=0$). As variáveis ficam ordenadas como:
\begin{equation}
\begin{array}{c|cccc}
\text{índice} & 0 & 1 & 2 & 3 \\
\hline
\text{variável} & x_{00} & x_{01} & x_{10} & x_{11}
\end{array}
\end{equation}
\listeq{Indexação linear das variáveis no exemplo}
A carga média é $L_{avg}=(0.3+0.5)/2=0.4$.

Antes de adicionar a penalidade one-hot, os termos diagonais do balanceamento são:
\begin{equation}
\begin{aligned}
Q_{0,0}^{bal}=Q_{1,1}^{bal}
&= w_0^2 - 2L_{avg}w_0
= 0.09 - 2(0.4)(0.3)
= -0.15,\\
Q_{2,2}^{bal}=Q_{3,3}^{bal}
&= w_1^2 - 2L_{avg}w_1
= 0.25 - 2(0.4)(0.5)
= -0.15.
\end{aligned}
\end{equation}
\listeq{Termos diagonais do balanceamento no exemplo}

Para o núcleo $0$, a expansão de $(L_0 - L_{avg})^2$ contém o termo de
co-localização:
\[
2w_0w_1x_{00}x_{10}=0.30x_{00}x_{10}.
\]
De forma análoga, a expansão de $(L_1 - L_{avg})^2$ contém:
\[
2w_0w_1x_{01}x_{11}=0.30x_{01}x_{11}.
\]
Como $Q$ é simétrica, este contributo é dividido pelas duas posições simétricas da matriz.

Assim, a matriz do objetivo de balanceamento é:
\begin{equation}
Q_{bal} =
\begin{pmatrix}
-0.15 & 0 & 0.15 & 0 \\
0 & -0.15 & 0 & 0.15 \\
0.15 & 0 & -0.15 & 0 \\
0 & 0.15 & 0 & -0.15
\end{pmatrix}.
\end{equation}
\listeq{Matriz de balanceamento do exemplo}

Com $P=1.0$, a penalidade one-hot subtrai $1.0$ a cada termo diagonal e adiciona $1.0$
entre as variáveis do mesmo processo em núcleos distintos: $x_{00}$ com $x_{01}$, e
$x_{10}$ com $x_{11}$. A matriz final fica:
\begin{equation}
Q =
\begin{pmatrix}
-1.15 & 1.0 & 0.15 & 0 \\
1.0 & -1.15 & 0 & 0.15 \\
0.15 & 0 & -1.15 & 1.0 \\
0 & 0.15 & 1.0 & -1.15
\end{pmatrix}.
\end{equation}
\listeq{Matriz Q final do exemplo}

Para verificar o efeito dos termos de co-localização, comparem-se duas soluções viáveis:
A, com os processos distribuídos ($x_{00}=x_{11}=1$), e B, com ambos os processos no
núcleo $0$ ($x_{00}=x_{10}=1$). Usando $E(\mathbf{x})=\mathbf{x}^{\top}Q\mathbf{x}$:
\begin{equation}
\begin{aligned}
E_A &= Q_{0,0} + Q_{3,3}
= -1.15 + (-1.15)
= \mathbf{-2.30},\\
E_B &= Q_{0,0} + Q_{2,2} + Q_{0,2} + Q_{2,0}
= -1.15 + (-1.15) + 0.15 + 0.15
= \mathbf{-2.00}.
\end{aligned}
\end{equation}
\listeq{Comparação de energia entre duas soluções viáveis}

Como $E_A < E_B$, a solução distribuída tem menor energia. A matriz penaliza, portanto,
a concentração dos dois processos no mesmo núcleo, favorecendo uma distribuição mais
equilibrada.

Depois de construída a matriz $Q$, o problema de escalonamento passa a estar expresso
numa forma compatível com otimizadores \acs{QUBO}. O próximo passo será resolver esta matriz, isto é, procurar a bitstring $\mathbf{x}$ que minimiza a energia $E(\mathbf{x})=\mathbf{x}^{\top}Q\mathbf{x}$. Neste trabalho, essa etapa é realizada
com \acs{QAOA}~\cite{Farhi2014QAOA}, implementado em PennyLane~\cite{Bergholm2018PennyLane}.
